\chapter{2018年全国I卷（理）}

\section{单选题}

\begin{problem}
设 \(z=\frac{1-\i}{1+\i}+2\i\)，则 \(\lvert z\rvert=\)（\quad）

\choices
    {\(0\)}
    {\(\frac12\)}
    {\(1\)}
    {\(\sqrt2\)}
\end{problem}

\begin{answer}
C
\end{answer}

\begin{solution}
\[
z=\frac{(1-\i)^2}{(1+\i)(1-\i)}+2\i
 =\frac{-2\i}{2}+2\i=\i.
\]
因此 \(\lvert z\rvert=1\).
\end{solution}

\begin{problem}
已知集合 \(A=\{x\mid x^2-x-2>0\}\)，则 \(C_{\R}A=\)（\quad）

\choices
    {\(\{x\mid -1<x<2\}\)}
    {\(\{x\mid -1\le x\le 2\}\)}
    {\(\{x\mid x<-1\}\cup\{x\mid x>2\}\)}
    {\(\{x\mid x\le-1\}\cup\{x\mid x\ge2\}\)}
\end{problem}

\begin{answer}
B
\end{answer}

\begin{solution}
因为
\[
x^2-x-2=(x+1)(x-2),
\]
所以
\[
A=(-\infty,-1)\cup(2,+\infty).
\]
因此
\[
C_{\R}A=[-1,2]=\{x\mid -1\le x\le2\}.
\]
故选 B.
\end{solution}

\begin{problem}
某地区经过一年的新农村建设，农村的经济收入增加了一倍，实现翻番. 为更好地了解该地区农村的经济收入变化情况，统计了该地区新农村建设前后农村的经济收入构成比例，得到如下饼图：

\examfiguregroup[float=inline]{img/2018/national_paper_1_science/q3_fig1.png}{%
    \bitmapinclude[max width=0.46\linewidth]{img/2018/national_paper_1_science/q3_fig1.png}\hfill
    \bitmapinclude[max width=0.46\linewidth]{img/2018/national_paper_1_science/q3_fig2.png}%
}

则下面结论中不正确的是（\quad）

\choices
    {新农村建设后，种植收入减少}
    {新农村建设后，其他收入增加了一倍以上}
    {新农村建设后，养殖收入增加了一倍}
    {新农村建设后，养殖收入与第三产业收入的总和超过了经济收入的一半}
\end{problem}

\begin{answer}
A
\end{answer}

\begin{solution}
设建设前农村经济收入总额为 \(T\)，则建设后总额为 \(2T\). 建设前种植、养殖、其他、第三产业收入分别为 \(0.60T\)，\(0.30T\)，\(0.04T\)，\(0.06T\)；建设后相应收入分别为 \(0.74T\)，\(0.60T\)，\(0.10T\)，\(0.56T\). 种植收入由 \(0.60T\) 增加到 \(0.74T\)，并未减少；其他收入为原来的 \(2.5\) 倍，养殖收入为原来的 \(2\) 倍，且建设后养殖收入与第三产业收入之和为 \(1.16T>T\). 因此不正确的是 A.
\end{solution}

\begin{problem}
记 \(S_n\) 为等差数列 \(\{a_n\}\) 的前 \(n\) 项和. 若 \(3S_3=S_2+S_4\)，\(a_1=2\)，则 \(a_5=\)（\quad）

\choices
    {\(-12\)}
    {\(-10\)}
    {\(10\)}
    {\(12\)}
\end{problem}

\begin{answer}
B
\end{answer}

\begin{solution}
设等差数列的公差为 \(d\)，则
\[
S_2=4+d,\qquad S_3=6+3d,\qquad S_4=8+6d.
\]
由题设
\[
3(6+3d)=(4+d)+(8+6d),
\]
解得 \(d=-3\). 因而
\[
a_5=a_1+4d=2+4(-3)=-10.
\]
故选 B.
\end{solution}

\begin{problem}
设函数 \(f(x)=x^3+(a-1)x^2+ax\). 若 \(f(x)\) 为奇函数，则曲线 \(y=f(x)\) 在点 \((0,0)\) 处的切线方程为（\quad）

\choices
    {\(y=-2x\)}
    {\(y=-x\)}
    {\(y=2x\)}
    {\(y=x\)}
\end{problem}

\begin{answer}
D
\end{answer}

\begin{solution}
由 \(f(x)\) 为奇函数，得
\[
-x^3+(a-1)x^2-ax=f(-x)=-f(x)=-x^3-(a-1)x^2-ax.
\]
比较 \(x^2\) 项系数，得 \(a-1=0\)，即 \(a=1\). 此时
\[
f(x)=x^3+x,
\qquad f'(x)=3x^2+1.
\]
又 \(f(0)=0\)，所以曲线在 \((0,0)\) 处的切线斜率为 \(f'(0)=1\)，切线方程为 \(y=x\). 故选 D.
\end{solution}

\begin{problem}
在 \(\triangle ABC\) 中，\(AD\) 为 \(BC\) 边上的中线，\(E\) 为 \(AD\) 的中点，则 \(\overrightarrow{EB}=\)（\quad）

\choices
    {\(\frac34\overrightarrow{AB}-\frac14\overrightarrow{AC}\)}
    {\(\frac14\overrightarrow{AB}-\frac34\overrightarrow{AC}\)}
    {\(\frac34\overrightarrow{AB}+\frac14\overrightarrow{AC}\)}
    {\(\frac14\overrightarrow{AB}+\frac34\overrightarrow{AC}\)}
\end{problem}

\begin{answer}
A
\end{answer}

\begin{solution}
设 \(\overrightarrow{AB}=\bs b\)，\(\overrightarrow{AC}=\bs c\). 因为 \(AD\) 是 \(BC\) 的中线，所以
\[
\overrightarrow{AD}=\frac{\bs b+\bs c}{2}.
\]
又因为 \(E\) 是 \(AD\) 的中点，故
\[
\overrightarrow{AE}=\frac{\bs b+\bs c}{4}.
\]
于是
\[
\overrightarrow{EB}=\overrightarrow{AB}-\overrightarrow{AE}
 =\frac34\bs b-\frac14\bs c
 =\frac34\overrightarrow{AB}-\frac14\overrightarrow{AC}.
\]
故选 A.
\end{solution}

\begin{problem}
某圆柱的高为 \(2\)，底面周长为 \(16\)，其三视图如图. 圆柱表面上的点 \(M\) 在正视图上的对应点为 \(A\)，圆柱表面上的点 \(N\) 在左视图上的对应点为 \(B\)，则在此圆柱侧面上，从 \(M\) 到 \(N\) 的路径中，最短路径的长度为（\quad）

\FigureLayoutDeclare{img/2018/national_paper_1_science/q7_fig1.png}{8.9cm}{10bp 32bp 0bp 0bp}
\bitmapfigure[max width=6.0cm]{img/2018/national_paper_1_science/q7_fig1.png}

\choices
    {\(2\sqrt{17}\)}
    {\(2\sqrt5\)}
    {\(3\)}
    {\(2\)}
\end{problem}

\begin{answer}
B
\end{answer}

\begin{solution}
将圆柱侧面沿一条母线展开，得到长为底面周长 \(16\)、宽为高 \(2\) 的矩形. 由三视图中 \(A\) 与 \(B\) 的位置可知，点 \(M\)，\(N\) 在圆柱底面圆周方向上相差四分之一周长，因此两点沿圆周方向的最短距离为
\[
\frac{16}{4}=4.
\]
展开后，\(M\)，\(N\) 的竖直方向距离为圆柱高 \(2\)，故最短路径为展开矩形中的直线段，其长度为
\[
\sqrt{2^2+4^2}=2\sqrt5.
\]
故选 B.
\end{solution}

\begin{problem}
设抛物线 \(C:y^2=4x\) 的焦点为 \(F\)，过点 \((-2,0)\) 且斜率为 \(\frac23\) 的直线与 \(C\) 交于 \(M\)，\(N\) 两点，则 \(\overrightarrow{FM}\cdot\overrightarrow{FN}=\)（\quad）

\choices
    {\(5\)}
    {\(6\)}
    {\(7\)}
    {\(8\)}
\end{problem}

\begin{answer}
D
\end{answer}

\begin{solution}
抛物线 \(y^2=4x\) 的焦点为 \(F=(1,0)\). 过 \((-2,0)\) 且斜率为 \(\dfrac23\) 的直线方程为
\[
y=\frac23(x+2).
\]
代入抛物线方程，得
\[
\frac49(x+2)^2=4x
\quad\Longleftrightarrow\quad
x^2-5x+4=0,
\]
所以交点的横坐标为 \(1\)，\(4\)，相应纵坐标为 \(2\)，\(4\). 不妨设 \(M=(1,2)\)，\(N=(4,4)\)，则
\[
\overrightarrow{FM}=(0,2),\qquad \overrightarrow{FN}=(3,4),
\]
从而
\[
\overrightarrow{FM}\cdot\overrightarrow{FN}=0\cdot3+2\cdot4=8.
\]
故选 D.
\end{solution}

\begin{problem}
已知函数 \(f(x)=\begin{cases}
\e^x,&x\le0,\\
\ln x,&x>0
\end{cases}\)，\(g(x)=f(x)+x+a\). 若 \(g(x)\) 存在 \(2\text{ 个}\) 零点，则 \(a\) 的取值范围是（\quad）
\choices
    {\([-1,0)\)}
    {\([0,+\infty)\)}
    {\([-1,+\infty)\)}
    {\([1,+\infty)\)}
\end{problem}

\begin{answer}
C
\end{answer}

\begin{solution}
函数 \(g\) 可写为
\[
g(x)=
\begin{cases}
\e^x+x+a,&x\le0,\\
\ln x+x+a,&x>0.
\end{cases}
\]
在 \((-\infty,0]\) 上，\(g'(x)=\e^x+1>0\)，且
\[
\lim_{x\to-\infty}g(x)=-\infty,\qquad g(0)=a+1.
\]
因此左侧分支有一个零点的充要条件是 \(a+1\ge0\)，即 \(a\ge-1\)，其中 \(a=-1\) 时零点为 \(x=0\). 在 \((0,+\infty)\) 上，\(g'(x)=\dfrac1x+1>0\)，且
\[
\lim_{x\to0^+}g(x)=-\infty,\qquad \lim_{x\to+\infty}g(x)=+\infty,
\]
所以右侧分支恒有且仅有一个零点. 故 \(g(x)\) 有两个零点的充要条件为 \(a\ge-1\)，即 \(a\in[-1,+\infty)\). 故选 C.
\end{solution}

\begin{problem}
如图来自古希腊数学家希波克拉底所研究的几何图形. 此图由三个半圆构成，三个半圆的直径分别为直角三角形 \(ABC\) 的斜边 \(BC\)，直角边 \(AB\)，\(AC\). \(\triangle ABC\) 的三边所围成的区域记为 \(I\)，黑色部分记为 \(II\)，其余部分记为 \(III\). 在整个图形中随机取一点，此点取自 \(I\)，\(II\)，\(III\) 的概率分别记为 \(p_1\)，\(p_2\)，\(p_3\)，则（\quad）

\bitmapfigure[max width=0.52\linewidth]{img/2018/national_paper_1_science/q10_fig1.png}

\choices
    {\(p_1=p_2\)}
    {\(p_1=p_3\)}
    {\(p_2=p_3\)}
    {\(p_1=p_2+p_3\)}
\end{problem}

\begin{answer}
A
\end{answer}

\begin{solution}
设以 \(BC\)，\(AB\)，\(AC\) 为直径的三个半圆面积分别为 \(S_{BC}\)，\(S_{AB}\)，\(S_{AC}\). 由勾股定理，\(BC^2=AB^2+AC^2\)，所以
\[
S_{BC}=S_{AB}+S_{AC}.
\]
由图形分区，直径为 \(BC\) 的半圆由区域 \(\mathrm{I}\)，\(\mathrm{III}\) 组成，而另外两个半圆合起来由区域 \(\mathrm{II}\)，\(\mathrm{III}\) 组成，因此
\[
S_{BC}=S_I+S_{III},\qquad S_{AB}+S_{AC}=S_{II}+S_{III}.
\]
两式结合得 \(S_I=S_{II}\). 在同一个整体图形中随机取点时，区域概率与区域面积成正比，故 \(p_1=p_2\). 故选 A.
\end{solution}

\begin{problem}
已知双曲线 \(C:\frac{x^2}{3}-y^2=1\)，\(O\) 为坐标原点，\(F\) 为 \(C\) 的右焦点，过 \(F\) 的直线与 \(C\) 的两条渐近线的交点分别为 \(M\)，\(N\). 若 \(\triangle OMN\) 为直角三角形，则 \(\lvert MN\rvert=\)（\quad）

\choices
    {\(\frac32\)}
    {\(3\)}
    {\(2\sqrt3\)}
    {\(4\)}
\end{problem}

\begin{answer}
B
\end{answer}

\begin{solution}
双曲线的两条渐近线为 \(y=\pm\dfrac{x}{\sqrt3}\)，右焦点为 \(F=(2,0)\). 设
\[
M=\left(u,\frac{u}{\sqrt3}\right),\qquad N=\left(v,-\frac{v}{\sqrt3}\right).
\]
由直线 \(MN\) 经过 \(F\)，利用横截距公式得
\[
\frac{2uv}{u+v}=2,
\qquad\text{即}\qquad uv=u+v.
\]
若三角形 \(OMN\) 在 \(O\) 处为直角，则
\[
\overrightarrow{OM}\cdot\overrightarrow{ON}
 =uv\left(1-\frac13\right)\ne0,
\]
不可能. 因此直角在 \(M\) 或 \(N\) 处.

当直角在 \(M\) 处时，
\[
\overrightarrow{MO}\cdot\overrightarrow{MN}=0
\quad\Longleftrightarrow\quad v=2u.
\]
代入 \(uv=u+v\)，得 \(2u^2=3u\). 因 \(u\ne0\)，所以 \(u=\dfrac32\)，\(v=3\). 当直角在 \(N\) 处时同理可得 \(u=3\)，\(v=\dfrac32\). 两种情形下
\[
\lvert MN\rvert^2=(u-v)^2+\frac{(u+v)^2}{3}
 =\left(\frac32\right)^2+\frac{\left(\frac92\right)^2}{3}=9.
\]
所以 \(\lvert MN\rvert=3\). 故选 B.
\end{solution}

\begin{problem}
已知正方体的棱长为 \(1\)，每条棱所在直线与平面 \(\alpha\) 所成的角都相等，则 \(\alpha\) 截此正方体所得截面面积的最大值为（\quad）

\choices
    {\(\frac{3\sqrt3}{4}\)}
    {\(\frac{2\sqrt3}{3}\)}
    {\(\frac{3\sqrt2}{4}\)}
    {\(\frac{\sqrt3}{2}\)}
\end{problem}

\begin{answer}
A
\end{answer}

\begin{solution}
取正方体的一组棱方向为坐标轴方向. 设平面 \(\alpha\) 的一个法向量为 \(\bs n=(n_1,n_2,n_3)\). 直线与平面的夹角相等，等价于
\[
\frac{|n_1|}{|\bs n|}=\frac{|n_2|}{|\bs n|}=\frac{|n_3|}{|\bs n|},
\]
所以平面法向量的三个分量绝对值相等. 通过对坐标轴方向作对称变换，不妨设平行平面的方程为
\[
x+y+z=s.
\]
令截面面积为 \(A(s)\). 当 \(0\le s\le1\) 时，截面是顶点为 \((s,0,0)\)，\((0,s,0)\)，\((0,0,s)\) 的等边三角形，故
\[
A(s)=\frac{\sqrt3}{2}s^2.
\]
当 \(1\le s\le2\) 时，可看作上述三角形去掉三个全等的小角三角形，因而
\[
A(s)=\frac{\sqrt3}{2}\bigl[s^2-3(s-1)^2\bigr].
\]
当 \(2\le s\le3\) 时，由正方体关于中心的对称性，\(A(s)=A(3-s)\). 在 \([1,2]\) 上，
\[
A'(s)=\frac{\sqrt3}{2}(-4s+6),
\]
所以截面面积在 \(s=\dfrac32\) 时取得最大值，且
\[
A\left(\frac32\right)=\frac{\sqrt3}{2}\left[\frac94-3\left(\frac12\right)^2\right]=\frac{3\sqrt3}{4}.
\]
故选 A.
\end{solution}

\section{填空题}

\begin{problem}
若 \(x\)，\(y\) 满足约束条件 \(\begin{cases}
x-2y-2\le0,\\
x-y+1\ge0,\\
y\le0
\end{cases}\)，则 \(z=3x+2y\) 的最大值为\fillinblank{}.
\end{problem}

\begin{answer}
6
\end{answer}

\begin{solution}
由约束条件 \(x-2y-2\le0\) 得 \(x\le2y+2\). 对固定的 \(y\)，因为目标函数 \(z=3x+2y\) 关于 \(x\) 增大，所以
\[
z\le3(2y+2)+2y=8y+6.
\]
又由 \(y\le0\)，得 \(z\le6\). 当 \(y=0\)，\(x=2\) 时满足全部约束，且 \(z=6\)，所以最大值为 \(6\).
\end{solution}

\begin{problem}
记 \(S_n\) 为数列 \(\{a_n\}\) 的前 \(n\) 项和. 若 \(S_n=2a_n+1\)，则 \(S_6=\)\fillinblank{}.
\end{problem}

\begin{answer}
-63
\end{answer}

\begin{solution}
当 \(n=1\) 时，\(S_1=a_1=2a_1+1\)，所以 \(a_1=-1\)，即 \(S_1=-1\). 对 \(n\ge2\)，由 \(a_n=S_n-S_{n-1}\) 及题设得
\[
S_n=2(S_n-S_{n-1})+1,
\]
从而
\[
S_n=2S_{n-1}-1.
\]
令 \(T_n=S_n-1\)，则 \(T_n=2T_{n-1}\)，且 \(T_1=-2\). 因此
\[
T_6=-2\cdot2^5=-64,
\qquad S_6=T_6+1=-63.
\]
\end{solution}

\begin{problem}
从 \(2\text{ 位}\) 女生，\(4\text{ 位}\) 男生中选 \(3\text{ 人}\) 参加科技比赛，且至少有 \(1\text{ 位}\) 女生入选，则不同的选法共有\fillinblank{}\(\text{种}\).
\end{problem}

\begin{answer}
16
\end{answer}

\begin{solution}
按入选女生人数分类：
\[
\mathrm C_2^1\mathrm C_4^2+\mathrm C_2^2\mathrm C_4^1
 =2\times6+1\times4=16.
\]
因此不同的选法共有 \(16\) 种.
\end{solution}

\begin{problem}
已知函数 \(f(x)=2\sin x+\sin2x\)，则 \(f(x)\) 的最小值是\fillinblank{}.
\end{problem}

\begin{answer}
\(-\dfrac{3\sqrt3}{2}\)
\end{answer}

\begin{solution}
令 \(c=\cos x\)，则
\[
f'(x)=2\cos x+2\cos2x
 =2\bigl(c+2c^2-1\bigr)=2(2c-1)(c+1).
\]
所以极值点满足 \(\cos x=\dfrac12\) 或 \(\cos x=-1\). 当 \(\cos x=\dfrac12\) 时，\(\sin x=\pm\dfrac{\sqrt3}{2}\)，且
\[
f(x)=2\sin x+2\sin x\cos x=3\sin x,
\]
其最小值为 \(-\dfrac{3\sqrt3}{2}\). 当 \(\cos x=-1\) 时，\(\sin x=0\)，函数值为 \(0\). 因此 \(f(x)\) 的最小值为
\[
-\frac{3\sqrt3}{2}.
\]
\end{solution}

\section{解答题}

\begin{problem}
在平面四边形 \(ABCD\) 中，\(\angle ADC=90^\circ\)，\(\angle A=45^\circ\)，\(AB=2\)，\(BD=5\).

\begin{enumerate}
\item 求 \(\cos\angle ADB\)；
\item 若 \(DC=2\sqrt2\)，求 \(BC\).
\end{enumerate}
\end{problem}

\begin{solution}
\begin{enumerate}
\item 由题意，\(\angle DAB=45^\circ\). 在 \(\triangle ABD\) 中，由正弦定理得
\[
\frac{AB}{\sin\angle ADB}=\frac{BD}{\sin\angle DAB}
\]
所以 \(\sin\angle ADB=\dfrac{2\sin45^\circ}{5}=\dfrac{\sqrt2}{5}\). 又因为 \(BD>AB\)，所以 \(\angle ADB<\angle DAB=45^\circ\)，从而
\[
\cos\angle ADB=\sqrt{1-\sin^2\angle ADB}=\sqrt{1-\frac{2}{25}}=\frac{\sqrt{23}}{5}.
\]

\item 因为 \(\angle ADC=90^\circ\)，所以 \(\angle BDC=90^\circ-\angle ADB\)，从而
\[
\cos\angle BDC=\sin\angle ADB=\frac{\sqrt2}{5}.
\]
在 \(\triangle BDC\) 中，由余弦定理得
\[
\begin{aligned}
BC^2&=BD^2+DC^2-2\cdot BD\cdot DC\cos\angle BDC\\
&=25+8-2\cdot5\cdot2\sqrt2\cdot\frac{\sqrt2}{5}=25.
\end{aligned}
\]
因为 \(BC>0\)，所以 \(BC=5\).
\end{enumerate}
\end{solution}

\begin{problem}
如图，四边形 \(ABCD\) 为正方形，\(E\)，\(F\) 分别为 \(AD\)，\(BC\) 的中点，以 \(DF\) 为折痕把 \(\triangle DFC\) 折起，使点 \(C\) 到达点 \(P\) 的位置，且 \(PF\perp BF\).

\begin{enumerate}
\item 证明：平面 \(PEF\perp\) 平面 \(ABFD\)；
\item 求 \(DP\) 与平面 \(ABFD\) 所成角的正弦值.
\end{enumerate}

\bitmapfigure[float=inline,width=0.55\linewidth]{img/2018/national_paper_1_science/q18_fig1.png}
\end{problem}

\begin{solution}
\begin{enumerate}
\item 因为四边形 \(ABCD\) 为正方形，且 \(E\)，\(F\) 分别为 \(AD\)，\(BC\) 的中点，所以 \(EF\parallel AB\)，而 \(AB\perp BF\)，从而 \(EF\perp BF\). 题设又给出 \(PF\perp BF\). 直线 \(PF\) 与 \(EF\) 相交于 \(F\)，且都在平面 \(PEF\) 内，所以 \(BF\perp\) 平面 \(PEF\). 又因为 \(BF\subset\) 平面 \(ABFD\)，故平面 \(PEF\perp\) 平面 \(ABFD\).

\item 取 \(BF=1\) 为长度单位，则正方形边长为 \(2\)，且 \(CF=1\). 折叠过程中点 \(D\)，\(F\) 不动，点 \(C\) 到达点 \(P\)，所以 \(DP=DC=2\)，\(PF=CF=1\).

由（1）知平面 \(PEF\perp\) 平面 \(ABFD\)，而 \(DE\subset\) 平面 \(ABFD\)，\(DE\perp EF\)，所以 \(DE\perp\) 平面 \(PEF\)，即 \(DE\perp PE\). 又 \(DE=\dfrac{AD}{2}=1\)，在直角三角形 \(DPE\) 中，
\[
PE=\sqrt{DP^2-DE^2}=\sqrt{2^2-1^2}=\sqrt3.
\]
由于 \(PE^2+PF^2=3+1=4=EF^2\)，所以 \(\angle EPF=90^\circ\). 设 \(PH\perp EF\)，垂足为 \(H\)，则 \(PH\) 是直角三角形 \(PEF\) 斜边 \(EF\) 上的高，故
\[
PH=\frac{PE\cdot PF}{EF}=\frac{\sqrt3}{2},\qquad EH=\frac{PE^2}{EF}=\frac32,\qquad HF=\frac{PF^2}{EF}=\frac12.
\]

以 \(H\) 为原点，取 \(\overrightarrow{FB}\) 方向为 \(x\) 轴正方向，取 \(\overrightarrow{HF}\) 方向为 \(y\) 轴正方向，取垂直于平面 \(ABFD\) 且指向 \(P\) 的方向为 \(z\) 轴正方向，则
\[
D\left(-1,-\frac32,0\right)
\]
且 \(P\left(0,0,\frac{\sqrt3}{2}\right)\). 因此
\[
\overrightarrow{DP}=\left(1,\frac32,\frac{\sqrt3}{2}\right)
\]
且 \(|DP|=2\).
平面 \(ABFD\) 的一个单位法向量为 \(\bs n=(0,0,1)\). 设 \(DP\) 与平面 \(ABFD\) 所成的角为 \(\theta\)，则
\[
\sin\theta=\frac{|\overrightarrow{DP}\cdot\bs n|}{|DP|}=\frac{\sqrt3/2}{2}=\frac{\sqrt3}{4}.
\]
\end{enumerate}
\end{solution}

\begin{problem}
设椭圆 \(C:\frac{x^2}{2}+y^2=1\) 的右焦点为 \(F\)，过 \(F\) 的直线 \(l\) 与 \(C\) 交于 \(A\)，\(B\) 两点，点 \(M\) 的坐标为 \((2,0)\).

\begin{enumerate}
\item 当 \(l\) 与 \(x\) 轴垂直时，求直线 \(AM\) 的方程；
\item 设 \(O\) 为坐标原点，证明：\(\angle OMA=\angle OMB\).
\end{enumerate}
\end{problem}

\begin{solution}
\begin{enumerate}
\item 椭圆的长半轴平方为 \(2\)，短半轴平方为 \(1\)，所以焦距的一半为 \(c=\sqrt{2-1}=1\)，从而右焦点为 \(F=(1,0)\). 当直线 \(l\) 与 \(x\) 轴垂直时，\(l\) 的方程为 \(x=1\). 代入椭圆方程得
\[
\frac12+y^2=1
\]
所以 \(A\)，\(B\) 为 \(\left(1,\dfrac{\sqrt2}{2}\right)\)，\(\left(1,-\dfrac{\sqrt2}{2}\right)\). 因此直线 \(AM\) 的方程有两种可能：
\[
y=-\frac{\sqrt2}{2}(x-2)\quad\text{或}\quad y=\frac{\sqrt2}{2}(x-2).
\]

\item 先考虑直线 \(l\) 为 \(x\) 轴的情形，此时 \(A\)，\(B\) 都在 \(x\) 轴上，所以 \(\angle OMA=\angle OMB=0^\circ\).

当 \(l\) 与 \(x\) 轴垂直时，由第（1）问可知，直线 \(AM\)，\(BM\) 的斜率绝对值相等，且 \(M\)，\(O\) 的连线为水平线，所以 \(\angle OMA=\angle OMB\).

下面设 \(l\) 既不与 \(x\) 轴重合，也不与其垂直，记其斜率为 \(k\)，则 \(k\ne0\)，且 \(l\) 的方程为 \(y=k(x-1)\). 设 \(A\)，\(B\) 的横坐标分别为 \(x_1\)，\(x_2\)，则 \(x_1\)，\(x_2\) 是方程
\[
\frac{x^2}{2}+k^2(x-1)^2=1
\]
的两个根，即
\[
(1+2k^2)x^2-4k^2x+2k^2-2=0.
\]
由韦达定理得
\[
x_1+x_2=\frac{4k^2}{1+2k^2},\qquad x_1x_2=\frac{2k^2-2}{1+2k^2}.
\]

因为椭圆上点的横坐标满足 \(|x_i|\le\sqrt2<2\)，所以直线 \(MA\)，\(MB\) 的斜率分别为
\[
m_1=\frac{k(x_1-1)}{x_1-2},\qquad m_2=\frac{k(x_2-1)}{x_2-2}.
\]
于是
\[
m_1+m_2=\frac{k\left[(x_1-1)(x_2-2)+(x_2-1)(x_1-2)\right]}{(x_1-2)(x_2-2)}.
\]
其中
\[
\begin{aligned}
&(x_1-1)(x_2-2)+(x_2-1)(x_1-2)\\
&=2x_1x_2-3(x_1+x_2)+4\\
&=\frac{4k^2-4-12k^2+4+8k^2}{1+2k^2}=0.
\end{aligned}
\]
故 \(m_1=-m_2\). 又因为 \(x_1<2\)，\(x_2<2\)，射线 \(MA\)，\(MB\) 都指向左侧，因而它们与射线 \(MO\) 所成角的正切值分别为 \(|m_1|\)，\(|m_2|\). 由 \(|m_1|=|m_2|\)，得 \(\angle OMA=\angle OMB\).
\end{enumerate}
\end{solution}

\begin{problem}
某工厂的某种产品成箱包装，每箱 \(200\text{ 件}\)，每一箱产品在交付用户之前要对产品作检验，如检验出不合格品，则更换为合格品. 检验时，先从这箱产品中任取 \(20\text{ 件}\) 作检验，再根据检验结果决定是否对余下的所有产品作检验. 设每件产品为不合格品的概率都为 \(p\)（\(0<p<1\)），且各件产品是否为不合格品相互独立.

\begin{enumerate}
\item 记 \(20\text{ 件}\) 产品中恰有 \(2\text{ 件}\) 不合格品的概率为 \(f(p)\)，求 \(f(p)\) 的最大值点 \(p_0\)；
\item 现对一箱产品检验了 \(20\text{ 件}\)，结果恰有 \(2\text{ 件}\) 不合格品，以（1）中确定的 \(p_0\) 作为 \(p\) 的值. 已知每件产品的检验费用为 \(2\text{ 元}\)，若有不合格品进入用户手中，则工厂要对每件不合格品支付 \(25\text{ 元}\) 的赔偿费用.
\begin{enumerate}
\item 若不对该箱余下的产品作检验，这一箱产品的检验费用与赔偿费用的和记为 \(X\)，求 \(EX\)；
\item 以检验费用与赔偿费用和的期望值为决策依据，是否该对这箱余下的所有产品作检验？
\end{enumerate}
\end{enumerate}
\end{problem}

\begin{solution}
\begin{enumerate}
\item 令前 \(20\text{ 件}\) 产品中不合格品的件数为随机变量，则恰有 \(2\text{ 件}\) 不合格品的概率为
\[
f(p)=\mathrm C_{20}^{2}p^2(1-p)^{18}.
\]
对 \(f(p)\) 求导，得
\[
f'(p)=\mathrm C_{20}^{2}p(1-p)^{17}\left[2(1-p)-18p\right]=2\mathrm C_{20}^{2}p(1-p)^{17}(1-10p).
\]
当 \(0<p<\dfrac1{10}\) 时，\(f'(p)>0\)；当 \(\dfrac1{10}<p<1\) 时，\(f'(p)<0\). 所以 \(f(p)\) 先增后减，最大值点为
\[
p_0=\frac1{10}.
\]

\item 由（1）知 \(p=p_0=\dfrac1{10}\).
\begin{enumerate}
\item 设余下的 \(180\text{ 件}\) 产品中不合格品的件数为随机变量 \(Y\)，则
\[
Y\sim B\left(180,\frac1{10}\right)
\]
且 \(EY=180\times\frac1{10}=18\).
若不检验余下产品，前 \(20\text{ 件}\) 的检验费用为 \(20\times2=40\text{ 元}\)，余下产品中每件不合格品产生 \(25\text{ 元}\) 的赔偿费用，因此
\[
X=40+25Y
\]
从而
\[
EX=40+25EY=40+25\times18=490\text{ 元}.
\]

\item 若检验余下的全部 \(180\text{ 件}\) 产品，则整箱产品共检验 \(200\text{ 件}\)，检验费用为
\[
200\times2=400\text{ 元}.
\]
检验出的不合格品会被更换为合格品，所以不会产生赔偿费用. 因为 \(400<490\)，所以应该对这箱产品余下的所有产品作检验.
\end{enumerate}
\end{enumerate}
\end{solution}

\begin{problem}
已知函数 \(f(x)=\frac1x-x+a\ln x\).

\begin{enumerate}
\item 讨论 \(f(x)\) 的单调性；
\item 若 \(f(x)\) 存在两个极值点 \(x_1\)，\(x_2\)，证明：\(\dfrac{f(x_1)-f(x_2)}{x_1-x_2}<a-2\).
\end{enumerate}
\end{problem}

\begin{solution}
\begin{enumerate}
\item 函数的定义域为 \((0,+\infty)\)，且
\[
f'(x)=-\frac1{x^2}-1+\frac ax=\frac{-x^2+ax-1}{x^2}=-\frac{x^2-ax+1}{x^2}.
\]

当 \(a\le2\) 时，\(x^2-ax+1\ge0\)，所以 \(f'(x)\le0\)，函数 \(f(x)\) 在 \((0,+\infty)\) 上单调递减. 当 \(a=2\) 时，\(f'(1)=0\)，但导数在其余点均为负，函数仍在整个定义域上单调递减.

当 \(a>2\) 时，令 \(f'(x)=0\)，得
\[
x^2-ax+1=0.
\]
记
\[
\alpha=\frac{a-\sqrt{a^2-4}}{2},\qquad \beta=\frac{a+\sqrt{a^2-4}}{2}.
\]
则 \(0<\alpha<\beta\)，且 \(\alpha\beta=1\). 因为
\[
-x^2+ax-1=-(x-\alpha)(x-\beta)
\]
所以 \(f(x)\) 在 \((0,\alpha)\) 上单调递减，在 \((\alpha,\beta)\) 上单调递增，在 \((\beta,+\infty)\) 上单调递减.

\item 函数有两个极值点时必有 \(a>2\). 不妨设两个极值点满足 \(0<x_1<x_2\)，由它们是方程 \(x^2-ax+1=0\) 的两个根，得
\[
x_1+x_2=a
\]
且 \(x_1x_2=1\).
特别地，令 \(t=x_2\)，则 \(t>1\)，且 \(x_1=\dfrac1t\). 由函数表达式得
\[
\begin{aligned}
\frac{f(x_1)-f(x_2)}{x_1-x_2}
&=\frac{\dfrac1{x_1}-\dfrac1{x_2}-(x_1-x_2)+a\ln\dfrac{x_1}{x_2}}{x_1-x_2}\\
&=-2+a\frac{\ln\dfrac{x_2}{x_1}}{x_2-x_1}\\
&=-2+a\frac{2\ln t}{t-\dfrac1t}.
\end{aligned}
\]
设 \(h(t)=t-\dfrac1t-2\ln t\)，则
\[
h'(t)=1+\frac1{t^2}-\frac2t=\frac{(t-1)^2}{t^2}>0\quad(t>1)
\]
且 \(h(1)=0\).
所以当 \(t>1\) 时，\(2\ln t<t-\dfrac1t\)，即
\[
\frac{2\ln t}{t-\dfrac1t}<1.
\]
又因为 \(a>2>0\)，从而
\[
\frac{f(x_1)-f(x_2)}{x_1-x_2}=-2+a\frac{2\ln t}{t-\dfrac1t}<a-2.
\]
\end{enumerate}
\end{solution}

\section{选考题：解答题}

\begin{problem}
在直角坐标系 \(xOy\) 中，曲线 \(C_1\) 的方程为 \(y=k\lvert x\rvert+2\). 以坐标原点为极点，\(x\) 轴正半轴为极轴建立极坐标系，曲线 \(C_2\) 的极坐标方程为 \(\rho^2+2\rho\cos\theta-3=0\).

\begin{enumerate}
\item 求 \(C_2\) 的直角坐标方程；
\item 若 \(C_1\) 与 \(C_2\) 有且仅有三个公共点，求 \(C_1\) 的方程.
\end{enumerate}
\end{problem}

\begin{solution}
\begin{enumerate}
\item 由极坐标与直角坐标的关系 \(x=\rho\cos\theta\)，\(y=\rho\sin\theta\)，\(\rho^2=x^2+y^2\)，原方程化为
\[
x^2+y^2+2x-3=0
\]
即
\[
(x+1)^2+y^2=4.
\]

\item 曲线 \(C_2\) 是圆心为 \(G=(-1,0)\)，半径为 \(2\) 的圆. 若 \(k\ge0\)，则 \(C_1\) 上的点满足 \(y\ge2\)，而圆 \(C_2\) 的最高点为 \((-1,2)\)，此时至多有一个公共点，不能满足题意. 因此必须有 \(k<0\).

当 \(k<0\) 时，先看 \(C_1\) 的左侧射线 \(x\le0\)，其方程为 \(y=-kx+2\). 联立圆的方程，得

\[
(1+k^2)x^2+(2-4k)x+1=0.
\]

由于
\[
\Delta=(2-4k)^2-4(1+k^2)=4k(3k-4)>0,
\]
且两根之和为 \(-\dfrac{2-4k}{1+k^2}<0\)，两根之积为 \(\dfrac1{1+k^2}>0\)，所以两根均为负数. 因此左侧射线与圆恰有两个公共点. 右侧射线在 \(x=0\) 与 \(x\to+\infty\) 时均在圆外，故一个交点对应相切、两个对应相交；总数为3时必相切.

当 \(x\ge0\) 时，\(C_1\) 的右侧射线为 \(y=kx+2\)，其所在直线可写为 \(kx-y+2=0\). 圆心到该直线的距离等于半径，故
\[
\frac{|2-k|}{\sqrt{1+k^2}}=2.
\]
因为 \(k<0\)，所以 \(2-k>0\)，从而
\[
(2-k)^2=4(1+k^2)
\]
解得 \(k=0\) 或 \(k=-\dfrac43\). 结合 \(k<0\)，得 \(k=-\dfrac43\). 此时切点为 \(\left(\dfrac35,\dfrac65\right)\)，确实在右侧射线上.

因此两条射线一共恰有三个公共点，故 \(C_1\) 的方程为
\[
y=-\frac43|x|+2.
\]
\end{enumerate}
\end{solution}

\begin{problem}
已知 \(f(x)=\lvert x+1\rvert-\lvert ax-1\rvert\).

\begin{enumerate}
\item 当 \(a=1\) 时，求不等式 \(f(x)>1\) 的解集；
\item 若 \(x\in(0,1)\) 时不等式 \(f(x)>x\) 成立，求 \(a\) 的取值范围.
\end{enumerate}
\end{problem}

\begin{solution}
\begin{enumerate}
\item 当 \(a=1\) 时，\(f(x)=|x+1|-|x-1|\). 分三种情况讨论：
\[
f(x)=
\begin{cases}
-2,&x<-1\\
2x,&-1\le x<1\\
2,&x\ge1
\end{cases}
\]
当 \(x<-1\) 时，\(f(x)=-2\)，不满足 \(f(x)>1\)；当 \(-1\le x<1\) 时，\(f(x)>1\) 等价于 \(2x>1\)，即 \(x>\dfrac12\)；当 \(x\ge1\) 时，\(f(x)=2>1\). 因此不等式的解集为
\[
\left(\frac12,+\infty\right).
\]

\item 因为 \(x\in(0,1)\)，所以 \(|x+1|=x+1\). 不等式 \(f(x)>x\) 等价于
\[
x+1-|ax-1|>x
\]
即
\[
|ax-1|<1.
\]
这又等价于
\[
-1<ax-1<1
\]
即
\[
0<ax<2.
\]
要使上式对一切 \(x\in(0,1)\) 成立，首先必须有 \(a>0\). 当 \(a>0\) 时，\(0<ax\) 对所有 \(x\in(0,1)\) 自动成立；而 \(ax<2\) 对所有 \(x<1\) 成立的充要条件是 \(a\le2\). 其中 \(a=2\) 时有 \(2x<2\)，仍对所有 \(x\in(0,1)\) 成立. 因此 \(a\) 的取值范围为
\[
(0,2].
\]
\end{enumerate}
\end{solution}
